\documentclass[12pt,aspectratio=43]{beamer}

% 自訂字體的封包
\usepackage{fontspec}

% 導入圖形與表格的封包
\usepackage{graphicx}
\usepackage{booktabs}

% 排列多個子圖形的封包
\usepackage{subfigure} 

 % 提供進階數學環境與公式排版
\usepackage{mathtools, amsmath, amsfonts, amsthm, latexsym} 

%% 設定中文字體
\usepackage{xeCJK}
\setCJKmainfont[AutoFakeBold=3.5 , AutoFakeSlant=0.2]{標楷體}
\setmainfont{Times New Roman}
\XeTeXlinebreaklocale "zh"
\XeTeXlinebreakskip = 0pt plus 1pt

% 自訂字體顏色的封包
\usepackage{color} 

% 圖表自動編號的封包
\usepackage{caption} 

% 允許表格的一格能多列呈現的封包
\usepackage{multirow} 

% 可指定表格排版的封包
\usepackage{array}

% 翻轉表格的封包
\usepackage{lscape} 

% 序列標號
\usepackage{enumerate} 

%% 設定自動編號
\setbeamertemplate{caption}[numbered]

%% 自訂顏色
\definecolor{Myblue}{RGB}{57,108,144}
\definecolor{Default_Blue}{RGB}{52,51,171}

% (動畫)字會若隱若現的顯示
\beamersetuncovermixins{\opaqueness<1->{45}}{\opaqueness<1->{95}} 

% 設定中文的標籤
\renewcommand{\figurename}{圖} 
\renewcommand{\tablename}{表} 

% 排版形式 
\mode<presentation>{
\usetheme{Berlin}
% enumerate 及 item 的形狀
\setbeamertemplate{itemize items}[triangle]
% 調整外框形式的字體大小
\setbeamerfont{headline}{size=\scriptsize}
\setbeamerfont{footline}{size=\scriptsize}
% 取消右下方的跳轉工具列
\setbeamertemplate{navigation symbols}{} 
% 自定義2：清除外框下緣但僅出頁碼
\setbeamertemplate{footline}[page number] 
% 顏色主題
\usecolortheme{dove}
% 設定頁面邊界
\setbeamersize{text margin left=0.6cm, text margin right=0.6cm}
\special{papersize=\the\paperwidth,\the\paperheight}
\providecommand{\tabularnewline}{\\}
}

\title{Meeting 1022}
\author{黃丞暘}
\institute{國立東華大學應數統計所}
\date{October 22, 2025}

\begin{document}

% 封面
\begin{frame}
  \titlepage
\end{frame}

% --- 目錄 ---
\begin{frame}{\textbf{目錄}}
  \tableofcontents
\end{frame}

%   資料介紹
\section{資料介紹}
\linespread{1}  %此行以後，皆為1倍行高
\begin{frame}[<+>]{\textbf{資料介紹}} % [<+>] 會自動以動畫呈現，就不用一直輸入 \pause。
\linespread{1.5} %此行以後，皆為1.5倍行高

    SECOM
	
	% 垂直距離的空格 (另外有\medskip, \smallskip, \vspace{?em})
    \bigskip 
	 
	 % 動畫出現至此的指令
	% \pause 
	
    \begin{itemize}[] % 封包 itemize 的 item
      \item 樣本數 1567
      \item 變數 Time／x1／x2／...／x590／PassFail
        \begin{enumerate}[] % 同一層：四個子項目並齊
          \item Features: 591, Y: PassFail (binary classification)
          \item Time: 2008-07-19 11:55:00
          \item PassFail: Pass $\to$ 0, Fail $\to$ 1 (14:1)
          \item Missing values: 41951
        \end{enumerate}
    \end{itemize}
\end{frame}


%	--- 文獻回顧I ---
\section{文獻回顧I}
\linespread{1}

% --- Frame 1: 圖片出現前的內容 ---
\begin{frame}[<+>]{\textbf{文獻回顧I}}
\linespread{1.5}
  \begin{itemize}
    \item Lai (2002)
    \begin{enumerate}
      \item 方法
      \begin{enumerate}
        \item Data preprocessing
        \item KNN-imputation
        \item Random forest features selection
        \item imbalance - SMOTE to 1:1
      \end{enumerate}
    \end{enumerate}
  \end{itemize}
\end{frame}

% --- Frame 2: 獨立顯示圖片 ---
\begin{frame}{\textbf{文獻回顧I-方法} (SMOTE)}
\centering
\begin{columns}[T,onlytextwidth]
  \column{0.5\textwidth}
  \centering
  \includegraphics[width=\linewidth]{C:/Users/USER/Documents/GitHub/psychic-spoon/0 Meeting/20251015_Meeting/SMOTE_BEFORE.png}\\
  {\scriptsize SMOTE before}

  \column{0.5\textwidth}
  \centering
  \includegraphics[width=\linewidth]{C:/Users/USER/Documents/GitHub/psychic-spoon/0 Meeting/20251015_Meeting/SMOTE_AFTER.png}\\
  {\scriptsize SMOTE after}
\end{columns}
\end{frame}

% --- Frame 3: 接續圖片之後的內容 ---
\begin{frame}{\textbf{文獻回顧I}}
\linespread{1.5}
  \begin{itemize}
    \item Lai (2002)
    \begin{enumerate}
      \item 方法
      \begin{enumerate}
        % --- 以下項目沒有 < > 標籤 ---
        % 這代表它們在這一頁出現時，會立即顯示 (通常是淡化狀態)
        % 作為動畫的背景/上下文
        \item Data preprocessing
        \item KNN-imputation
        \item Random forest features selection
        \item imbalance - SMOTE to 1:1

        % --- 動畫從這裡開始 ---
        % 第一個 <+-> 標籤，代表按下「下一步」後，才會出現的第一個新項目
        \item<+-> 0.8 training, 0.2 testing by decision tree
      \end{enumerate}
      % 後續的動畫會依序播放
      \item<+-> 結果
      \begin{enumerate}
        \item<+-> Best: ACC 86.94\%, FAR 9.22\%
        \item<+-> 反查 Fail 節點
      \end{enumerate}
    \end{enumerate}
  \end{itemize}
\end{frame}


% --- 研究動機 ---
\section{研究動機}

\linespread{1}
\begin{frame}{\textbf{研究動機}}
\linespread{1.5}

	\begin{itemize}[]
		\item K-folds cross validation
		\begin{enumerate}[]
      \item Split to training\&testing first    
      \item Robust performance
		\end{enumerate}
	\end{itemize}

	\begin{itemize}[]
		\item Metric
		\begin{enumerate}[]
			\item ACC, FAR, Sensitivity, Specificity
		\end{enumerate}
	\end{itemize}
    
  \begin{itemize}[]
		\item $y = \hat{f}(x)$
		\begin{enumerate}[]
			\item Important features
		\end{enumerate}
	\end{itemize}
\end{frame}

% --- 研究結果 ---
\section{研究結果}

\linespread{1}
\begin{frame}{\textbf{研究結果}( Decision Tree（Test）$k=223$ )}
\linespread{1.0}

\begin{columns}[T,onlytextwidth]

  % ===== 左欄：數據與混淆矩陣 =====
  \column{0.48\textwidth}

  \begin{block}
  \small
  \begin{tabular}{@{}ll@{}}
    Accuracy    & 0.8662 \\
    FAR         & 0.0922 \\
    Sensitivity & 0.2857 \\
    Specificity & 0.9078 \\
    GM          & 0.5093 \\
  \end{tabular}
  \end{block}

  \begin{block}{Confusion Matrix}
  \small
  \setlength{\tabcolsep}{5pt}
  \renewcommand{\arraystretch}{1.15}
  \begin{tabular}{c|cc}
           & \textbf{Pred 0} & \textbf{Pred 1} \\ \hline
  \textbf{True 0} & 266 & 27 \\
  \textbf{True 1} & 15  & 6  \\
  \end{tabular}
  \end{block}

  % ===== 右欄：圖 =====
  \column{0.52\textwidth}
  \centering
  \includegraphics[width=\linewidth]{C:/Users/USER/Documents/GitHub/psychic-spoon/0 Meeting/20251015_Meeting/validACC.png}
  % 若需要圖編號與標題，解除下面註解（需 capt-of）
  % \captionof{figure}{Validation accuracy across $k$}

\end{columns}

\end{frame}

\linespread{1}
\begin{frame}{\textbf{研究結果}}
\linespread{1.5}
	
	\begin{figure}
		\centering % 置中
		\includegraphics[scale=0.3]{C:/Users/USER/Documents/GitHub/psychic-spoon/0 Meeting/20251015_Meeting/PCA.png}\\ % 導入圖片並調整大小
        % 記得調整\parbox的寬度
        \parbox{0.9\textwidth}{\scriptsize} % 調整位置並輸入資料來源
		% \caption{} % 圖片名稱
	\end{figure}

\end{frame}

% --- 參考資料 ---
\section{參考資料}
\linespread{1}  %此行以後，皆為1倍行高
\begin{frame}{\textbf{參考資料}}
\linespread{1.5} %此行以後，皆為1.5倍行高

\footnotesize

\begin{thebibliography}{99} 

	\bibitem[Li-Fu Lai (2022)]{中文文獻範例} % 引注時顯示的名稱
		Li-Fu Lai (2022)。 % 作者 (年份)
		\newblock 使用KNN和決策樹方法對填補缺失值的製造數據進行失效預測和原因追蹤分析 % 文章名稱
		
\end{thebibliography}

\end{frame}

% --- 結束 ---
\section*{結束}

\begin{frame}[plain] % 產生無外框的頁面
	\begin{center}
		{\Huge The End}
		\bigskip\bigskip 
		
		{\LARGE Questions  \&  Comments}
	\end{center}
\end{frame}
\end{document}ㄒ
